\documentclass[11pt]{article}
\usepackage[margin=1in]{geometry}
\usepackage{enumitem}
% =============================================================================
% Preambles/header.tex — shared LaTeX/Beamer preamble for this project
%
% Use from a Beamer lecture by prepending:
%     \input{header}
% and compiling with TEXINPUTS=../Preambles:$TEXINPUTS (the /compile-latex
% skill does this automatically).
%
% PALETTE CONTRACT. Color names defined here MUST match the names in
% Quarto/theme-template.scss so Beamer and Quarto renderings use the same
% palette. The scripts/check-palette-sync.sh script enforces this.
%
% To customize: edit the HEX values below AND the matching $variable in
% Quarto/theme-template.scss, then run scripts/check-palette-sync.sh.
% =============================================================================

% -----------------------------------------------------------------------------
% Engine detection + encoding
%
% Our /compile-latex skill uses XeLaTeX, where inputenc is unnecessary and
% can produce encoding warnings. Only load inputenc under pdfTeX.
% -----------------------------------------------------------------------------
\usepackage{iftex}
\ifPDFTeX
  \usepackage[utf8]{inputenc}
\fi

\usepackage{amsmath, amssymb, amsthm}
\usepackage{booktabs}
\usepackage{graphicx}

% -----------------------------------------------------------------------------
% PALETTE — mirrors Quarto/theme-template.scss variable names.
% Load xcolor and define colors BEFORE anything that references them
% (notably hyperref, hypersetup, and the Beamer color assignments).
% -----------------------------------------------------------------------------
\usepackage{xcolor}

\definecolor{primary-blue}{HTML}{012169}      % headings, accents
\definecolor{primary-gold}{HTML}{B9975B}      % emphasis, borders
\definecolor{highlight-yellow}{HTML}{F2A900}  % markers, alerts
\definecolor{light-bg}{HTML}{E8EDF5}          % title / section backgrounds
\definecolor{jet}{HTML}{1A1A1A}               % body text

% Semantic colors (used in diagrams and annotations)
\definecolor{positive}{HTML}{15803D}          % good / observed / identified
\definecolor{negative}{HTML}{B91C1C}          % bad / problematic / bias
\definecolor{neutral}{HTML}{525252}           % reference / context

% Highlight accents (match .hi-* classes in the SCSS)
\definecolor{hi-slate}{HTML}{314F4F}
\definecolor{hi-green}{HTML}{2E8B57}
\definecolor{hi-red}{HTML}{C41E3A}

% Semantic aliases so skill templates and snippets can write
% \textcolor{accent}{...} without hard-coding "primary-blue".
\colorlet{accent}{primary-blue}
\colorlet{accent2}{primary-gold}

% -----------------------------------------------------------------------------
% Hyperref — loaded AFTER xcolor + palette so link colors resolve.
% -----------------------------------------------------------------------------
\usepackage{hyperref}
\hypersetup{colorlinks=true, linkcolor=primary-blue, urlcolor=primary-blue,
            citecolor=primary-blue, pdfborder={0 0 0}}

% -----------------------------------------------------------------------------
% Course code — set once per deck (after \input{header}, before \begin{document})
% via \coursecode{CS401}. Shown in the footer on every slide so a deck's course
% is identifiable without opening the title page. Empty by default.
% -----------------------------------------------------------------------------
\makeatletter
\newcommand{\coursecode}[1]{\gdef\@coursecodetext{#1}}
\gdef\@coursecodetext{}
\makeatother

% -----------------------------------------------------------------------------
% Beamer theme assignments (only apply under Beamer).
%
% We detect Beamer via \@ifclassloaded{beamer}, which is the documented,
% non-internal way. This must be wrapped in \makeatletter/\makeatother
% because \@ifclassloaded contains an @-catcode character.
% -----------------------------------------------------------------------------
\makeatletter
\@ifclassloaded{beamer}{%
  \usefonttheme{professionalfonts}
  \setbeamercolor{structure}{fg=primary-blue}
  \setbeamercolor{frametitle}{fg=primary-blue}
  \setbeamercolor{title}{fg=primary-blue}
  \setbeamercolor{subtitle}{fg=primary-gold}
  \setbeamercolor{author}{fg=primary-blue}
  \setbeamercolor{institute}{fg=neutral}
  \setbeamercolor{date}{fg=primary-gold}
  \setbeamercolor{itemize item}{fg=primary-blue}
  \setbeamercolor{itemize subitem}{fg=primary-gold}
  \setbeamercolor{alerted text}{fg=primary-gold}
  \setbeamercolor{block title}{fg=white, bg=primary-blue}
  \setbeamercolor{block body}{bg=light-bg}
  % Semantic block variants: block=definition/concept (blue), exampleblock=
  % worked example (green/positive), alertblock=key takeaway or caution (gold).
  % Keep to <=2 colored boxes per slide across ALL three variants (INV-7).
  \setbeamercolor{block title example}{fg=white, bg=positive}
  \setbeamercolor{block body example}{bg=light-bg}
  \setbeamercolor{block title alerted}{fg=white, bg=primary-gold}
  \setbeamercolor{block body alerted}{bg=light-bg}

  % Minimal footer: course code (left) + page X / N (right), no navigation clutter
  \setbeamertemplate{navigation symbols}{}
  \setbeamertemplate{footline}{%
    \color{neutral}\footnotesize\hspace{0.7em}\@coursecodetext\hfill
    \insertframenumber/\inserttotalframenumber\hspace{0.7em}\vspace{0.3em}}
}{}
\makeatother

% -----------------------------------------------------------------------------
% TikZ setup — libraries the snippet gallery uses
% -----------------------------------------------------------------------------
\usepackage{tikz}
\usetikzlibrary{arrows.meta, positioning, calc, decorations.pathreplacing,
                fit, shapes.geometric, backgrounds}

% Shared TikZ styles — snippets and hand-written diagrams can pick these up
% via `\begin{tikzpicture}[every node/.style={...}]` or by naming specific
% styles (e.g., `\node[dag-node] (X) at (0,0) {X};`).
\tikzset{
  % Node styles for causal DAGs and similar
  dag-node/.style = {
    draw=primary-blue, fill=light-bg, rounded corners=2pt,
    minimum width=1.2cm, minimum height=0.9cm, align=center, font=\small
  },
  decision-node/.style = {
    draw=primary-gold, fill=white, diamond, aspect=1.8,
    minimum width=1.8cm, minimum height=1.2cm, align=center, font=\small,
    inner sep=1pt
  },
  % Edge styles — observed vs counterfactual
  observed-edge/.style = {-{Latex[length=2.5mm]}, thick, draw=jet},
  counterfactual-edge/.style = {-{Latex[length=2.5mm]}, thick, draw=neutral, dashed},
  confound-edge/.style = {-{Latex[length=2.5mm]}, thick, draw=negative, dashed},
  % Data-point styles for plots
  observed-dot/.style = {fill=primary-blue, draw=primary-blue, circle, inner sep=1.2pt},
  counterfactual-dot/.style = {fill=white, draw=neutral, circle, inner sep=1.2pt}
}

% -----------------------------------------------------------------------------
% Convenience macros
% -----------------------------------------------------------------------------
\newcommand{\muted}[1]{\textcolor{neutral}{#1}}
\newcommand{\key}[1]{\textcolor{primary-gold}{\textbf{#1}}}
\newcommand{\good}[1]{\textcolor{positive}{#1}}
\newcommand{\bad}[1]{\textcolor{negative}{#1}}

% Standout frame macro: full-bleed dark background for section transitions.
% Beamer-only (uses \begin{frame}/\setbeamercolor/\usebeamerfont) -- do not
% call from an article-class file (e.g. Notes/<CODE>/*-notes.tex). \muted,
% \key, \good, \bad, and \coursecode above are plain xcolor/gdef and are
% safe to use from either class; verified when Notes/article-class support
% was added.
% Expand: \transitionslide{Your transition title}
\newcommand{\transitionslide}[1]{%
  \begingroup
  \setbeamercolor{background canvas}{bg=primary-blue}
  \begin{frame}[plain]
    \centering
    \vfill
    {\usebeamerfont{title}\color{white}\Large #1\par}
    \vfill
  \end{frame}
  \endgroup
}

% Section divider frame (Beamer-only), an alternative to \transitionslide with a
% darker two-line style: near-black (jet) background, a small white label line
% above a larger gold title, and a thin gold rule along the bottom edge.
% Expand: \sectiondivider{Section 2}{Pointers}
\newcommand{\sectiondivider}[2]{%
  \begingroup
  \setbeamercolor{background canvas}{bg=jet}
  \begin{frame}[plain]
    \vspace*{3.0cm}
    {\color{white}\ttfamily\large #1\par}
    \vspace{0.5cm}
    {\color{primary-gold}\LARGE #2\par}
    \begin{tikzpicture}[remember picture, overlay]
      \draw[primary-gold, line width=2.5pt]
        ($(current page.south west)+(0,0.1)$) -- ($(current page.south east)+(0,0.1)$);
    \end{tikzpicture}
  \end{frame}
  \endgroup
}

% -----------------------------------------------------------------------------
% Channel watermark — "Concepts That Click" (YouTube).
%
% Article-class files (CompetitiveExam Q/A sets, Notes, Assignments) get a
% light diagonal draft watermark on every page plus a centered footer naming
% the channel. Beamer decks get a subtle TikZ background watermark instead
% (the footline already carries the course code). Centralized here so every
% MCQ PDF is branded without per-file edits.
% -----------------------------------------------------------------------------
\makeatletter
\@ifclassloaded{article}{%
  \usepackage{draftwatermark}
  \SetWatermarkText{Concepts That Click}
  \SetWatermarkScale{1}
  \SetWatermarkFontSize{2cm}
  \SetWatermarkLightness{0.86}
  \SetWatermarkAngle{45}
  \usepackage{fancyhdr}
  \pagestyle{fancy}
  \fancyhf{}
  \fancyfoot[C]{\footnotesize\textcolor{neutral}{Concepts That Click~|~YouTube: \href{https://www.youtube.com/@concepts.thatclick}{@concepts.thatclick}}}
  \renewcommand{\headrulewidth}{0pt}
  \renewcommand{\footrulewidth}{0pt}
  \fancypagestyle{plain}{%
    \fancyhf{}%
    \fancyfoot[C]{\footnotesize\textcolor{neutral}{Concepts That Click~|~YouTube: \href{https://www.youtube.com/@concepts.thatclick}{@concepts.thatclick}}}%
    \renewcommand{\headrulewidth}{0pt}%
    \renewcommand{\footrulewidth}{0pt}%
  }
}{}
\@ifclassloaded{beamer}{%
  \setbeamertemplate{background}{%
    \begin{tikzpicture}[remember picture, overlay]
      \node[rotate=45, opacity=0.055, align=center]
        at (current page.center)
        {\Huge\textcolor{neutral}{Concepts That Click}};
    \end{tikzpicture}%
  }
}{}
\makeatother 
\coursecode{JKSSB}
\renewcommand{\thesection}{61.\arabic{section}}
\newtheorem{example}{Example}[section]
\newenvironment{definitionbox}[1]{%
  \par\noindent\textbf{Definition (#1).}\ \itshape
}{\par\vspace{0.3em}}
\title{Current Office, Email and Statement Evaluation --- Detailed Lecture Notes}
\author{JKSSB Computer Awareness}
\date{\today}

\begin{document}
\maketitle

\section{Word: structured content and content controls}
Most Word documents are free-form, but many office documents --- forms,
templates, and records --- need \emph{structure}: a fixed place for a name, a
date, a choice from a list, or a check box. Word provides \textbf{content
controls} for exactly this task. A content control is an individual element
that you insert at a chosen point in a document or template; the user can type
or choose only what the control allows.

\begin{definitionbox}{Content control}
A single, structured element inserted into a Word document or template to
capture one piece of information in a controlled way.
\end{definitionbox}

Common content controls include plain text, rich text, picture, drop-down
list, combo box, date picker, and check box. They are inserted from the
\textbf{Developer} tab, which is hidden by default and is enabled through
Word Options $\rightarrow$ Customize Ribbon. A content control can restrict
editing at that spot and can enforce limited \textbf{data validation}: a
drop-down list allows only the listed items, and a date picker allows only a
date. This is the fact behind the JA 2026 word item (PYQ-12): a content
control \emph{can} enforce data validation, so the statement that it cannot is
the incorrect one.

Content controls should not be confused with \emph{placeholders}, which belong
to PowerPoint slide layouts, or with form \emph{fields} from older Word
versions. In this course, ``content control'' means the Word Developer-tab
element that structures a document.

\section{First-line and hanging indents}
Indentation decides how far a paragraph's text sits from the left (or right)
margin. Word distinguishes two special paragraph indents.

\begin{center}
\begin{tabular}{p{0.22\textwidth}p{0.66\textwidth}}
\toprule
\textbf{Indent} & \textbf{Meaning} \\
\midrule
First-line indent & Only the \textbf{first line} of the paragraph is indented;
the remaining lines start at the left margin. \\
Hanging indent & Every line \textbf{except the first} is indented; the first
line hangs out to the left. Used for references, bibliographies, and bulleted
lists. \\
\bottomrule
\end{tabular}
\end{center}

\begin{definitionbox}{Hanging indent}
A paragraph indent in which all lines except the first are pushed inward, so
the first line ``hangs'' at the left margin.
\end{definitionbox}

These are paragraph-level formatting settings. Two paragraphs can carry
different indents even when they use the \emph{same style}, because a
paragraph style supplies a default that individual paragraphs can override.
The usual wrong statement claims that a hanging indent shifts the first line;
in fact the first line is the one that stays out, and all the others move
inward.

\section{The Navigation Pane}
The \textbf{Navigation Pane} is a Word panel opened from the \textbf{View}
tab. It is a browsing aid rather than a formatting tool. It presents three
views or tabs.

\begin{enumerate}
  \item \textbf{Headings}: an outline of the document built from paragraphs
    that carry a \emph{Heading style} (Heading 1, Heading 2, and so on).
  \item \textbf{Pages}: thumbnail pictures of each page.
  \item \textbf{Results}: the matches for text typed into the search box.
\end{enumerate}

Because the Headings view is driven by heading styles, applying a Heading
style does two things at once: it gives the paragraph an outline level and it
makes the paragraph appear in the Navigation Pane (and in an automatic Table
of Contents).

\begin{example}
A student writes a long report. Applying \emph{Heading 1} to each chapter
title makes every chapter appear in the Navigation Pane, so the reader can
jump straight to any chapter. A reviewer's \emph{comment} on a sentence does
\emph{not} appear in the Navigation Pane; comments are read in the
markup/comments pane. That distinction is the trap in JA 2026 (PYQ-12).
\end{example}

\section{Track Changes and revision metadata}
\textbf{Track Changes} is Word's review tool. When it is switched on, every
edit is recorded as a \emph{revision} instead of being applied silently. Three
kinds of change are tracked: insertions, deletions, and formatting changes.
Each recorded change stores \textbf{metadata} --- the \textbf{author's name}
and the \textbf{date and time} of the edit.

\begin{definitionbox}{Revision metadata with Track Changes}
The information Word stores alongside a tracked change: who made it and when.
\end{definitionbox}

A tracked change can be \textbf{accepted} or \textbf{rejected}. Accepting a
change keeps the new text as final and \textbf{removes the revision metadata}
attached to that change; rejecting a change restores the original text. Before
sharing a document, Word's \emph{Inspect Document} feature can be used to
strip hidden personal information such as author names.

\section{Section Break versus Page Break}
Word documents are divided into \textbf{sections}, and page-level settings
belong to a section. The two breaks behave quite differently.

\begin{center}
\begin{tabular}{p{0.18\textwidth}p{0.36\textwidth}p{0.36\textwidth}}
\toprule
\textbf{Break} & \textbf{What it does} & \textbf{What it changes} \\
\midrule
Page Break & Moves the following text to the next page. & Nothing else; the
document stays in the same section with the same headers, footers, and
margins. \\
Section Break & Starts a \textbf{new section}. & Each section may have its own
headers and footers, margins, page orientation, columns, and page numbering. \\
\bottomrule
\end{tabular}
\end{center}

\begin{definitionbox}{Section Break}
A break that ends the current section and begins a new one, so that
page-level formatting can differ from that point onward.
\end{definitionbox}

\begin{example}
A report needs no header on the title page but a running header on every later
page. Inserting a single Page Break will not do it, because both pages remain
in one section. The correct tool is a \textbf{Section Break} after the title
page, after which the new section can carry its own header. This is the exact
distinction tested in JA 2026 (PYQ-13): Section Break, not Page Break, allows
different headers and footers in one document.
\end{example}

\section{Excel: nested IF, blank and non-numeric input}
An \textbf{IF} function chooses between two values:
\texttt{=IF(condition, value\_if\_true, value\_if\_false)}. A \textbf{nested
IF} places a second IF inside one of those two positions, which lets one
formula test several conditions in turn:
\[
\texttt{=IF(A1>50,"High",IF(A1>25,"Medium","Low"))}.
\]
Excel evaluates the conditions from left to right and returns the value of the
\textbf{first} condition that is TRUE; once a condition is satisfied, the
remaining tests are skipped.

\begin{example}
With A1 = 70 the first condition (70 > 50) is TRUE, so the result is
\emph{High}. With A1 = 30 the first condition is FALSE but the second
(30 > 25) is TRUE, so the result is \emph{Medium}. With A1 = 10 both are
FALSE, so the result is \emph{Low}.
\end{example}

The behaviour on \textbf{blank} and \textbf{non-numeric} input is the trap
(PYQ-14). In a numeric comparison Excel treats a \textbf{blank cell as zero}.
Therefore, if A1 is empty, \texttt{A1>50} is FALSE and
\texttt{A1>25} is also FALSE, so the formula falls through to the final value
--- \emph{Low}. A blank cell does not make the formula stop or produce an
error; it is simply read as 0.

Non-numeric \textbf{text} is different again. Text is not a number, so:
\begin{itemize}
  \item arithmetic on it fails --- for example \texttt{=A1+5} with text in A1
    returns the error \texttt{\#VALUE!};
  \item \texttt{ISNUMBER} returns FALSE for a text cell and for a blank cell.
\end{itemize}
For this reason, robust formulas guard the input, for instance with
\texttt{ISBLANK}, \texttt{ISNUMBER}, or \texttt{IFERROR}, before using the
value in a calculation.

\section{COUNTIFS and SUMPRODUCT}
Counting rows that meet \emph{several} conditions at once is a frequent task.
Excel offers two routes.

\begin{definitionbox}{COUNTIFS}
A function that counts the rows satisfying \emph{all} of the given
range-and-criteria pairs (AND logic):
\texttt{COUNTIFS(range1, criteria1, range2, criteria2, \ldots)}.
\end{definitionbox}

The \textbf{COUNTIFS} function takes \emph{paired} arguments: a range followed
by the criterion that applies to it, repeated for each condition. All the
ranges must be the same size. The singular \textbf{COUNTIF} handles one
condition; the plural \textbf{COUNTIFS} handles several. Text criteria can use
the wildcards \texttt{*} (any sequence of characters) and \texttt{?} (any one
character).

\begin{definitionbox}{SUMPRODUCT}
A function that multiplies corresponding elements of the given arrays and
returns the sum of those products.
\end{definitionbox}

On its own, \texttt{SUMPRODUCT(A1:A3,B1:B3)} returns
\texttt{A1*B1 + A2*B2 + A3*B3}. It can also be used to \emph{count}, by
building boolean arrays and multiplying them:

\begin{center}
\texttt{=SUMPRODUCT((A1:A10>50)*(B1:B10="Yes"))}
\end{center}

\noindent counts the rows in which the first condition AND the second
condition are both true, because the product is 1 only when both parts are 1.
The double-unary form \texttt{=SUMPRODUCT(--(A1:A10>50))} is a common way to
coerce TRUE/FALSE into 1/0.

\begin{example}
To count how many students scored above 50 \emph{and} belong to the Science
stream, either
\texttt{=COUNTIFS(Marks,">50",Stream,"Science")} or
\texttt{=SUMPRODUCT((Marks>50)*(Stream="Science"))} will work. The stable
distinction tested in JA 2026 (PYQ-15, TRAP-14) is that COUNTIFS takes several
criteria directly, while SUMPRODUCT does arithmetic on arrays and can
therefore both \emph{add} and \emph{count} --- it does not ``only add.''
\end{example}

\section{Email: filters, signatures, POP3 and privacy}
\subsection*{Filters and rules}
An email \textbf{filter}, usually called a \textbf{rule} in desktop mail
clients, acts automatically on messages. A rule has two parts: a
\textbf{condition} (sender, recipient, subject, keywords, presence of an
attachment) and an \textbf{action} (move to a folder, flag, forward, delete,
mark as read). A typical rule sends all mail from one mailing list into a
dedicated folder, keeping the inbox tidy.

\subsection*{Digital signature versus encryption}
These two security features are constantly confused (PYQ-16, TRAP-15).

\begin{center}
\begin{tabular}{p{0.22\textwidth}p{0.34\textwidth}p{0.34\textwidth}}
\toprule
\textbf{Point} & \textbf{Digital signature} & \textbf{Encryption} \\
\midrule
Purpose & Authenticity and integrity & Confidentiality \\
What it establishes & Who sent the message, and that it was not altered in
transit & That only the intended reader can read the contents \\
Does it hide the content? & No & Yes \\
\bottomrule
\end{tabular}
\end{center}

A \textbf{digital signature} does \emph{not} encrypt the message body. It
proves the sender's identity and detects tampering. It is \textbf{encryption}
that scrambles the content so that only the holder of the correct key can read
it.

\subsection*{POP3 and IMAP}
\textbf{POP3} stands for \textbf{Post Office Protocol version 3}. It is a mail
\emph{retrieval} protocol: it \textbf{downloads} messages from the mail server
to the local device. By default it typically \textbf{removes} the messages
from the server, so once downloaded the mail lives only on that machine and is
not kept in sync across devices. \textbf{IMAP} works differently: it leaves
the mail on the server and keeps all devices synchronised. The stable
distinction (PYQ-16, TRAP-17) is that POP3 downloads to the local machine,
whereas IMAP syncs. \textbf{SMTP}, in contrast, is the protocol for
\emph{sending} mail.

\subsection*{BCC and recipient privacy}
An email has three recipient fields.

\begin{center}
\begin{tabular}{p{0.14\textwidth}p{0.22\textwidth}p{0.52\textwidth}}
\toprule
\textbf{Field} & \textbf{Full form} & \textbf{Visibility} \\
\midrule
To & --- & Primary recipients; visible to everyone. \\
Cc & Carbon Copy & Copied recipients; visible to everyone. \\
Bcc & Blind Carbon Copy & Addresses hidden from the \emph{other} recipients. \\
\bottomrule
\end{tabular}
\end{center}

The key point (PYQ-16, TRAP-16) is what Bcc hides. Bcc hides the blind
recipients from \emph{each other} and from the To/Cc recipients; it does
\emph{not} hide them from the \textbf{sender}, who always sees the complete
recipient list. A common wrong statement says that Bcc hides the sender and
recipients from one another --- the sender is never hidden.

\section{PowerPoint SmartArt conversion}
\textbf{SmartArt} turns a plain list into a diagram. PowerPoint can convert
content in both directions, with one firm limit.

\begin{itemize}
  \item \textbf{Text to SmartArt}: select a bulleted list and choose
    \textbf{Home} $\rightarrow$ Convert to SmartArt Graphic; choose a layout.
    The text is preserved as the diagram's content.
  \item \textbf{SmartArt to text}: with the graphic selected, use SmartArt
    Tools Design $\rightarrow$ Convert $\rightarrow$ Convert to Text. You get
    a bulleted list back; the shapes and any pictures in the graphic are
    \emph{lost}.
  \item \textbf{SmartArt to shapes}: Convert to Shapes breaks the graphic
    into ordinary shapes. Once that is done, the shapes \textbf{cannot} be
    converted back into SmartArt.
\end{itemize}

\begin{example}
A presenter has six bullet points and wants a cycle diagram. She selects the
list and uses \emph{Convert to SmartArt}. Later she needs the plain list
again, so she uses \emph{Convert to Text} and the six points return as
bullets. What cannot be undone is the step to individual shapes: after
\emph{Convert to Shapes}, rebuilding the SmartArt would mean starting over
(the ``re-entry'' point behind PYQ-17).
\end{example}

The direct source of Convert to SmartArt is \emph{text}. A picture or a chart
is not transformed into SmartArt in one step; you would instead insert a
picture-type SmartArt layout and place the images by hand.

\section{Exam retrieval and common confusions}
The following distinctions prevent most errors in this unit.
\begin{enumerate}
  \item Different headers and footers require a \textbf{Section Break}, not a
    Page Break (TRAP-13).
  \item A \textbf{hanging indent} shifts every line except the first; a
    first-line indent shifts only the first line.
  \item A \textbf{blank cell} is treated as 0 in a comparison; it does not
    cause an error (PYQ-14).
  \item \textbf{COUNTIFS} counts with AND logic over paired arguments;
    \textbf{SUMPRODUCT} multiplies arrays and can both add and count
    (TRAP-14).
  \item A \textbf{digital signature} gives authenticity and integrity, not
    confidentiality; \textbf{encryption} gives confidentiality (TRAP-15).
  \item \textbf{POP3 downloads} to the local machine; \textbf{IMAP} syncs
    (TRAP-17).
  \item \textbf{Bcc} hides recipients from each other, never from the sender
    (TRAP-16).
  \item \textbf{Navigation Pane} shows headings, pages, and search results,
    not comments.
\end{enumerate}

\begin{example}
An item states, ``A digital signature encrypts the message so that only the
recipient can read it.'' This is wrong. The signature establishes the sender
and integrity; encryption is the feature that hides the content. The wrong
option has merged two distinct security services.
\end{example}

\subsection*{Exercises}
\begin{enumerate}[label=\arabic*.]
  \item Distinguish a first-line indent from a hanging indent, with one use
    for each.
  \item State why a Section Break, and not a Page Break, is needed for
    different headers and footers.
  \item Explain what a nested IF returns when several conditions are true at
    once.
  \item What does \texttt{=IF(A1>50,"High","Low")} display when A1 is blank,
    and why?
  \item Give the general form of COUNTIFS, and state one thing SUMPRODUCT can
    do that COUNTIFS cannot.
  \item Distinguish a digital signature from encryption in one sentence.
  \item What does POP3 do with mail, and how does that differ from IMAP?
  \item In a message, who can see the addresses placed in the Bcc field?
  \item Describe two directions of SmartArt conversion and one that cannot be
    reversed.
\end{enumerate}

\subsection*{Solutions}
\begingroup\small
\begin{enumerate}[label=\arabic*.]
  \item A first-line indent indents only the first line (normal paragraph
    openings); a hanging indent indents every line except the first
    (references and bulleted lists).
  \item Page-level settings such as headers and footers belong to a section.
    A Page Break keeps the same section, so its header/footer cannot differ;
    a Section Break starts a new section that can carry its own header and
    footer.
  \item A nested IF returns the value of the \emph{first} condition that is
    TRUE; once a condition is satisfied the remaining tests are skipped.
  \item It displays \emph{Low}, because a blank cell is treated as 0, and 0 is
    not greater than 50, so the FALSE branch is taken.
  \item \texttt{COUNTIFS(range1, criteria1, range2, criteria2, \ldots)}.
    SUMPRODUCT can multiply corresponding array elements and return their sum
    (arithmetic), so it can add as well as count; COUNTIFS only counts.
  \item A digital signature proves the sender's identity and that the message
    was not altered, whereas encryption hides the content so only the
    intended reader can read it.
  \item POP3 downloads mail from the server to the local device and by
    default removes it from the server; IMAP leaves mail on the server and
    keeps devices synchronised.
  \item The sender can see all Bcc addresses, and each Bcc recipient is hidden
    from the other recipients (To, Cc, and the other Bcc recipients).
  \item Text converts to SmartArt (Convert to SmartArt), and SmartArt
    converts back to text (Convert to Text). Converting SmartArt to shapes is
    \emph{not} reversible --- the shapes cannot be turned back into SmartArt.
\end{enumerate}
\endgroup
\end{document}
